%!TEX TS-program = xelatex
%!TEX encoding = UTF-8 Unicode
% Awesome CV LaTeX Template for CV/Resume
%
% This template has been downloaded from:
% https://github.com/posquit0/Awesome-CV
%
% Author:
% Claud D. Park <posquit0.bj@gmail.com>
% http://www.posquit0.com
%
% Template license:
% CC BY-SA 4.0 (https://creativecommons.org/licenses/by-sa/4.0/)
%


%-------------------------------------------------------------------------------
% CONFIGURATIONS
%-------------------------------------------------------------------------------
% A4 paper size by default, use 'letterpaper' for US letter
\documentclass[11pt, a4paper]{awesome-cv}

% * Both of these packages are for getting page numbers in the footer.
\usepackage{lastpage}
\usepackage{fancyhdr}

% Configure page margins with geometry
\geometry{left=1.2cm, top=.8cm, right=1.2cm, bottom=1.6cm, footskip=.5cm}

% Color for highlights
% Awesome Colors: awesome-emerald, awesome-skyblue, awesome-red, awesome-pink, awesome-orange
%                 awesome-nephritis, awesome-concrete, awesome-darknight
\colorlet{awesome}{awesome-skyblue}
% Uncomment if you would like to specify your own color
% \definecolor{awesome}{HTML}{3E6D9C}

% Colors for text
% Uncomment if you would like to specify your own color
% \definecolor{darktext}{HTML}{414141}
% \definecolor{text}{HTML}{333333}
% \definecolor{graytext}{HTML}{5D5D5D}
% \definecolor{lighttext}{HTML}{999999}
% \definecolor{sectiondivider}{HTML}{5D5D5D}

% Set false if you don't want to highlight section with awesome color
\setbool{acvSectionColorHighlight}{true}

% If you would like to change the social information separator from a pipe (|) to something else
\renewcommand{\acvHeaderSocialSep}{\quad\textbar\quad}


%-------------------------------------------------------------------------------
%	PERSONAL INFORMATION
%	Comment any of the lines below if they are not required
%-------------------------------------------------------------------------------
% Available options: circle|rectangle,edge/noedge,left/right
% \photo[rectangle,edge,right]{./examples/profile}
\name{Udit}{Desai}
\position{Principal Software Engineer}
\address{Sydney, NSW, Australia}

\mobile{+61 459 343 918}
\email{desaiuditd@gmail.com}
\homepage{blog.incognitech.in}
\github{desaiuditd}
\linkedin{desaiuditd}
% \gitlab{gitlab-id}
% \stackoverflow{SO-id}{SO-name}
% \twitter{@twit}
% \x{x-id}
% \skype{skype-id}
% \reddit{reddit-id}
% \medium{madium-id}
% \kaggle{kaggle-id}
% \hackerrank{hackerrank-id}
% \telegram{telegram-username}
% \googlescholar{googlescholar-id}{name-to-display}
%% \firstname and \lastname will be used
% \googlescholar{googlescholar-id}{}
% \extrainfo{extra information}

\quote{``Walk the talk. Turn over every stone. Keep it human.''}


%-------------------------------------------------------------------------------
\begin{document}

% Print the header with above personal information
% Give optional argument to change alignment(C: center, L: left, R: right)
\makecvheader[C]

% Print the footer with 3 arguments(<left>, <center>, <right>)
% Leave any of these blank if they are not needed
\makecvfooter
  {\today}
  {Udit Desai~~~·~~~Résumé}
  % ! '*' after pageref is to avoid the link in last page number.
  % Ref: https://stackoverflow.com/a/4236418
  % Ref: https://stackoverflow.com/questions/4236389/page-x-of-y-in-latex-footer#comment122395019_4236418
  {\thepage\hspace{0pt} / \pageref*{LastPage}}


%-------------------------------------------------------------------------------
%	CV/RESUME CONTENT
%	Each section is imported separately, open each file in turn to modify content
%-------------------------------------------------------------------------------
%-------------------------------------------------------------------------------
%	SECTION TITLE
%-------------------------------------------------------------------------------
\cvsection{Summary}


%-------------------------------------------------------------------------------
%	CONTENT
%-------------------------------------------------------------------------------
\begin{cvparagraph}
Principal Software Engineer with 10+ years building publishing platforms at scale.
Building the tools that empower newsrooms behind
\href{https://www.smh.com.au/}{\color{awesome}Sydney Morning Herald},
\href{https://www.theage.com.au/}{\color{awesome}The Age} \&
\href{https://www.afr.com/}{\color{awesome}Australian Financial Review}
to deliver quality journalism efficiently.

Currently, shaping multi-team technical direction for AI agent engineering at
Nine --- building a custom GitHub Copilot plugin that encodes Nine's own engineering
standards, helping teams ship from idea to production without cutting corners.

Architecting Nine's Publishing Platform --- connecting a React block editor to
backend APIs (Go/GraphQL/Kafka) via a custom 2-way transform layer. Replaced a legacy
PHP monolith to unlock cross-team contribution and reduce dependency on a single
backend language.

Earned the Tech CTO Award (``Turn Over Every Stone'') for persistence in solving hard
problems. I also mentor engineers, help them grow in their careers, and build consensus
across engineering and executive stakeholders.
\end{cvparagraph}

\input{resume-template-parts/02-skills.tex}
\input{resume-template-parts/03-experience.tex}
%-------------------------------------------------------------------------------
%	SECTION TITLE
%-------------------------------------------------------------------------------
\cvsection{Technical Leadership}


%-------------------------------------------------------------------------------
%	CONTENT
%-------------------------------------------------------------------------------
\begin{cventries}

%---------------------------------------------------------
  \cventry
    {Novel CMS Architecture --- Gutenberg + Microservices Bridge} % Role/Title
    {Nine Publishing} % Organization
    {Sydney, Australia} % Location
    {2024 -- Present} % Date(s)
    {
      \begin{cvitems}
        \item {Problem: A WordPress monolith blocked cross-team contribution. Every code change required PHP knowledge, and newsroom editors needed better tools. But we couldn't move fast enough on legacy codebase.}
        \item {Solution: Re-build of the CMS layer to connect a React block editor to our internal APIs (Go/GraphQL/Kafka) --- bypassing WordPress entirely for content storage.}
            \begin{itemize}
                \item {2-way transform layer between Gutenberg blocks and internal ``Ink Blocks.''}
                    \begin{itemize}
                        \item {Write path: Gutenberg block array $\rightarrow$ Ink Blocks $\rightarrow$ GraphQL $\rightarrow$ gRPC $\rightarrow$ Kafka.}
                        \item {Read path: reverse. Required reverse-engineering Gutenberg's internal data flow.}
                    \end{itemize}
                \item {Reused select \texttt{@wordpress/} packages (\texttt{@wordpress/blocks}, \texttt{@wordpress/data}, \texttt{@wordpress/rich-text}) with custom equivalents for the rest.}
                \item {gRPC with proto contracts for compile-time schema enforcement. Kafka decouples authoring from downstream publication.}
                \item {Built PoC first, then showed it to leadership. Near production launch.}
            \end{itemize}
        \item {Result: Increased development velocity. Other engineering teams now contribute using React/TypeScript. Reduced reliance on PHP.}
      \end{cvitems}
    }

%---------------------------------------------------------
  \cventry
    {Engineering Excellence --- AI, Tooling \& Capability Building} % Role/Title
    {Nine Publishing} % Organization
    {Sydney, Australia} % Location
    {2023 -- Present} % Date(s)
    {
      \begin{cvitems}
        \item {Built AI agents for Git commit message generation and unit test writing}
            \begin{itemize}
                \item {Picked these two because they're high-impact and low-risk for AI augmentation.}
            \end{itemize}
        \item {Introduced Orbstack (Docker Desktop replacement), RenovateBot for dependency updates, and post-deploy Playwright smoke tests across the team.}
        \item {Coaching \& Mentoring: Engineers on the team grew to handle complex frontend work independently.}
            \begin{itemize}
                \item {Built a culture of knowledge sharing and continuous learning.}
                \item {Leadership style: open-ended questions over directives; explain the ``why''; show vulnerability to lower authority barriers; rubber-ducking with screen share.}
            \end{itemize}
        \item {Ran Engineers' Round Table forum and Deep Dive technical sessions. Established patterns and conventions that multiplied team velocity.}
      \end{cvitems}
    }

%---------------------------------------------------------
\end{cventries}

%---------------------------------------------------------
  \cventry
    {Brightcove Decommission -- Problem $\rightarrow$ Advocacy $\rightarrow$ PoC $\rightarrow$ Influence $\rightarrow$ Execution $\rightarrow$ Quality $\rightarrow$ Launch} % Role/Title
    {Nine Publishing} % Organization
    {Sydney, Australia} % Location
    {2026} % Date(s)
    {
      \begin{cvitems}
        \item {Problem: The legacy Brightcove plugin kept causing production incidents --- external vendor API failures were taking down the app.}
        \item {Solution: Built a replacement using React components already in the new architecture. Pitched the PoC to leadership.}
        \item {Execution: Offered to lead the rollout myself since the team was busy with other priorities.}
        \item {Quality: Conducted bug bash, wrote test cases document to help the team, fixed critical bugs before launch.}
        \item {Result: Zero production incidents from this source since launch.}
      \end{cvitems}
    }

\smallskip

\input{resume-template-parts/05-notable-projects.tex}
%-------------------------------------------------------------------------------
%	SECTION TITLE
%-------------------------------------------------------------------------------
\cvsection{Awards \& Recognition}


%-------------------------------------------------------------------------------
%	SUBSECTION TITLE
%-------------------------------------------------------------------------------
\cvsubsection{Company Awards}


%-------------------------------------------------------------------------------
%	CONTENT
%-------------------------------------------------------------------------------
\begin{cvhonors}

%---------------------------------------------------------
  \cvhonor
    {Tech CTO Award --- ``Turn Over Every Stone''} % Award
    {for persistence in solving hard problems, Nine Publishing} % Event
    {Sydney, Australia} % Location
    {Aug 2023} % Date(s)

%---------------------------------------------------------
  \cvhonor
    {INMA Award --- Ink Authoring} % Award
    {\href{https://www.inma.org/blogs/value-content/post.cfm/transparency-collaboration-drive-in-house-planning-innovations-at-nine}{\color{awesome}International News Media Association}} % Event
    {} % Location
    {Nov 2023} % Date(s)

%---------------------------------------------------------
  \cvhonor
    {Ink AI SEO Headlines ---Outstanding Team Work} % Award
    {Nine Hackathon} % Event
    {Sydney, Australia} % Location
    {Jul 2023} % Date(s)

%---------------------------------------------------------
\end{cvhonors}

%-------------------------------------------------------------------------------
%	SUBSECTION TITLE
%-------------------------------------------------------------------------------
\cvsubsection{Hackathon Awards}


%-------------------------------------------------------------------------------
%	CONTENT
%-------------------------------------------------------------------------------
\begin{cvhonors}

%---------------------------------------------------------
  \cvhonor
    {Best use of AWS} % Award
    {BroncoHack, \href{https://github.com/desaiuditd/slycam}{\color{awesome}Github}} % Event
    {Santa Clara, USA} % Location
    {Apr 2017} % Date(s)

%---------------------------------------------------------
  \cvhonor
    {Best use of Microsoft Azure} % Award
    {HackUCSC, \href{https://github.com/desaiuditd/outtatheblues}{\color{awesome}Github}} % Event
    {UC Santa Cruz, USA} % Location
    {Jan 2017} % Date(s)

%---------------------------------------------------------
  \cvhonor
    {Top 6 Finalists} % Award
    {HackRice, \href{https://github.com/desaiuditd/wayfarer}{\color{awesome}Github}} % Event
    {Rice University, USA} % Location
    {Nov 2016} % Date(s)

%---------------------------------------------------------
  \cvhonor
    {Best WordPress Data Hack} % Award
    {\href{https://techcrunch.com/video/incognitech-2/}{\color{awesome}TechCrunch Disrupt}, \href{https://github.com/desaiuditd/incognitech}{\color{awesome}Github}} % Event
    {San Francisco, USA} % Location
    {Sep 2016} % Date(s)

%---------------------------------------------------------
  \cvhonor
    {Best User Voice Experience} % Award
    {HackDavis, \href{https://github.com/desaiuditd/roomscore}{\color{awesome}Github}} % Event
    {UC Davis, USA} % Location
    {Jun 2016} % Date(s)

%---------------------------------------------------------
\end{cvhonors}

%-------------------------------------------------------------------------------
%	SECTION TITLE
%-------------------------------------------------------------------------------
\cvsection{Other Work}


%-------------------------------------------------------------------------------
%	SUBSECTION TITLE
%-------------------------------------------------------------------------------
\cvsubsection{Open Source}


%-------------------------------------------------------------------------------
%	CONTENT
%-------------------------------------------------------------------------------
\begin{cventries}

%---------------------------------------------------------
  \cventry
    {WordPress Contributor} % Role/Title
    {WordPress / Gutenberg} % Organization
    {} % Location
    {2014 -- 2021} % Date(s)
    {
      \begin{cvitems}
        \item {Contributions in \href{https://github.com/WordPress/gutenberg/pulls?q=is:pr+author:desaiuditd}{\color{awesome}Gutenberg}}
        \item {\href{https://core.trac.wordpress.org/report/25?USER=desaiuditd}{\color{awesome}Contributions} in \href{https://github.com/WordPress/wordpress-develop/pulls?q=is:pr+author:desaiuditd}{\color{awesome}WordPress Core}}
        \item {Other \href{https://profiles.wordpress.org/desaiuditd/\#content-plugins}{\color{awesome}WordPress Plugins}}
      \end{cvitems}
    }

%---------------------------------------------------------
\end{cventries}


%-------------------------------------------------------------------------------
%	SUBSECTION TITLE
%-------------------------------------------------------------------------------
\cvsubsection{Freelance Work}


%-------------------------------------------------------------------------------
%	CONTENT
%-------------------------------------------------------------------------------
\begin{cventries}

%---------------------------------------------------------
  \cventry
    {Web Engineer} % Role/Title
    {10up} % Organization
    {Remote / EMEA (UK)} % Location
    {Jun 2022 -- Sep 2023} % Date(s)
    {}

%---------------------------------------------------------
  \cventry
    {Lead Engineer} % Role/Title
    {Moxie Media Group} % Organization
    {Remote / USA} % Location
    {Dec 2014 -- Aug 2015} % Date(s)
    {}

%---------------------------------------------------------
  \cventry
    {WordPress Engineer} % Role/Title
    {XWP.co Pty. Ltd.} % Organization
    {Remote / Australia} % Location
    {Feb 2014 -- Aug 2015} % Date(s)
    {}

%---------------------------------------------------------
\end{cventries}

%-------------------------------------------------------------------------------
%	SECTION TITLE
%-------------------------------------------------------------------------------
\cvsection{Education}


%-------------------------------------------------------------------------------
%	CONTENT
%-------------------------------------------------------------------------------
\begin{cventries}

%---------------------------------------------------------
  \cventry
    {M.S. Computer Science} % Degree
    {Santa Clara University} % Institution
    {Santa Clara, CA} % Location
    {June 2017} % Date(s)
    {}

%---------------------------------------------------------
  \cventry
    {M.S. Information Technology} % Degree
    {DA-IICT} % Institution
    {Gandhinagar, India} % Location
    {May 2013} % Date(s)
    {}

%---------------------------------------------------------
  \cventry
    {B.S. Computer Science} % Degree
    {University of Pune} % Institution
    {Pune, India} % Location
    {June 2011} % Date(s)
    {}

%---------------------------------------------------------
\end{cventries}



%-------------------------------------------------------------------------------
\end{document}
